\section{Affine realization: gauges, elimination and witnesses}
\label{app:realization-bookkeeping}

\subsection{Coordinate choices and unused rows}
The nonzero vectors of $H=\Ftwo^2$ are the three states. Every permutation of them
extends uniquely to an element of $\operatorname{GL}(2,2)\cong S_3$. Applying this
map simultaneously to all state constants, successor differences and occurrence
expressions gives an isomorphic profile system. Swapping the two local paired
roles replaces $z_v$ by $z_v+1$; changing an anchor replaces comparisons by their
sums. Thus the construction does not depend on these presentational choices.
An independent permutation at one mask, without transforming the successor
links to other masks, is not a symmetry. Nor can different occurrences of one
absolute mask be independently rotated by grid symmetries.

At a load-two vertex the unused state is $b(a_M)+a_M(z_v+1)$. It is unused at that
vertex, not necessarily at the other vertices of mask $M$. When both paired input
states occur somewhere at $M$, $\kappa_M$ is their uniquely determined successor
difference. When only one occurs, it is an auxiliary description of the observed
transitions; it does not prescribe an unobserved output. In particular, the
unobserved row must not be assigned an affine extrapolation that equals the zero
vector. The arbitrary legal completion in the sufficiency proof avoids this
error. Endpoint comparisons cannot be dropped even when both endpoints have
load one: identical endpoint masks and states refer to the same controller row.

\subsection{Exact completion by elimination}
\label{app:linear-completion}
Fix one profile and partition its variables into retained variables $x$ and
variables $y$ to be completed. Write the system as
\[
 A_xx+A_yy=b.
\]
For a fixed $x$ there exists a completion precisely when
\[
 \lambda^T A_y=0\quad\Longrightarrow\quad
 \lambda^T A_xx=\lambda^T b.
\]
Indeed, this is exactly the condition that $b-A_xx$ belong to the column space
of $A_y$, by Gaussian elimination. This is the correct completion statement for
load-two roles: additional constraints on retained roles can remain after
elimination. Before fixing a profile the set of completions is a union of these
affine sets and need not itself be affine. An automatic completion assertion
based only on saturated correspondences is therefore not used.

\subsection{Signed graphs and incidence pseudoforests}
\label{app:forest-cases}
In a fixed profile, first combine repeated terms in each scalar equation and
retain constant equations. If every remaining equation has at most two variables,
form a signed multigraph with a reference vertex $*$ of value zero. Represent
$z_u=b$ by an edge $u*$ of gain $b$, and $z_u+z_v=b$ by an edge $uv$ of gain $b$.
Parallel edges are retained. A constant $0=1$ rejects the system.

\begin{proposition}[Graph balance]\label{prop:affine-graph-balance}
After constant equations are checked, such a system is consistent if and only
if the XOR of the gains around every cycle is zero. In particular, a forest is
consistent, and a component with one cycle has exactly its one parity condition.
\end{proposition}
\begin{proof}
A solution telescopes around each cycle. Conversely, choose a root in each
component and propagate values along a spanning tree, fixing the root of the
component containing $*$ to zero. Each non-tree edge is satisfied precisely when
its fundamental cycle has gain zero. This also proves that the solution space
has dimension $c-1$, where $c$ is the number of components including the one
containing $*$.
\end{proof}

The incidence graph of a general scalar system is instead bipartite, with variable
vertices and equation vertices. Its being a forest is \emph{not} by itself a
consistency criterion: the two equations $z=0$ and $z=1$ have an incidence tree.
Use the following exact reductions. Check constant equations; substitute every
unary fixation; discard degree-zero variables as free; and remove a degree-one
variable together with its unique incident equation, saving that equation for
back-substitution. Each reduction preserves existence of a solution. A finite
incidence forest has a leaf, so these reductions decide it completely.
If each incidence component has at most one cycle, the residual nonempty core
has degree two at every vertex and has equations
\[
 z_1+z_2=c_1,\ \ldots,\ z_{k-1}+z_k=c_{k-1},\ z_k+z_1=c_k.
\]
It is consistent precisely when $\sum_i c_i=0$. The constants here are the
\emph{updated} constants after substitution, not necessarily the original ones.
Solving the core and back-substituting reconstructs a solution. All these tests
must use one shared profile across the components.

For the general system an inclusion-minimal inconsistent set of rows is a
minimal binary dependency whose coefficients cancel and whose right-hand sides
sum to one. To see the minimality assertion, a proper coefficient dependency
with right-hand side one would itself contradict minimality; one with right-hand
side zero could be removed, leaving a smaller contradiction. These row circuits,
not necessarily graph cycles, are the general algebraic obstructions. A failed
profile is not a failed specification: every admissible profile must fail.

\subsection{The saturated NAND witness}
\label{app:nand-witnesses}
For the specification of Proposition~\ref{prop:nand8}, the occurrence word is
\[
(0,1,0,1,2,3,2,3,4,5,6,7,6,7,6,5,4,3,2,1).
\]
With ordinary state names $0,1,2$, the following assignments, in this occurrence
order, realize the three allowed pairs of saturated correspondence bits:
\begin{align*}
00:\quad &(0,0,2,1,1,0,2,1,1,1,1,0,2,1,0,2,0,2,0,2),\\
01:\quad &(0,0,2,1,1,0,2,1,1,0,2,1,1,0,0,2,2,2,0,2),\\
10:\quad &(0,0,2,1,1,2,0,1,1,1,1,2,0,1,2,2,0,0,2,2).
\end{align*}
The $x$-bit compares the W-pairs at occurrences $(1,19)$ and $(5,17)$;
the $y$-bit compares the E-pairs at $(4,6)$ and $(10,12)$. Reading the successors
verifies that each repeated $(q,M)$ has the same direction and next state.
States at each spatial vertex are distinct. These explicit assignments establish
the positive cases without an enumeration argument; the main-text four-state
contradiction excludes $11$.

\subsection{The six-vertex staircase}
\label{app:staircase}
Take directions $\E\N\E\N\E$, doubled edge indices $1,3$, and phases $0,1$,
respectively. Its word is
\[
 (0,1,2,1,2,3,4,5,4,3,4,3,2,1).
\]
All four internal vertices are saturated. Put $w=q_{13}=q_{11}$ for the
singleton W-state of mask $\N\W$, and $e=q_4=q_6$ for the singleton E-state
of mask $\E\Sdir$. Let $x$ compare the N-pairs $(q_1,q_3)$ and $(q_5,q_9)$,
and let $y$ compare the S-pairs $(q_2,q_{12})$ and $(q_8,q_{10})$.
The N-rows have successor pairs $(q_2,e)$ and $(e,q_{10})$.
Identity correspondence would force $q_2=e$, contradicting the two different
direction roles at one mask. Hence $x=1$, and the swapped correspondence forces
$q_2=q_{10}$, so $y=1$. The S-rows have successor pairs $(q_3,w)$ and $(q_9,w)$;
$y=1$ would force $q_3=w$, again contradicting distinct roles. Thus $y=0$.
Equivalently these constraints contain
\[
 x=1,\qquad x+y=0,\qquad y=0,
\]
whose sum is $0=1$. This proves impossibility directly. The separate assertion
that no smaller direction-consistent specification is impossible is an exact
finite check, not part of this argument.
